\section{Experimental Setup}
\label{sec:setup}

\paragraph{The suite.} LIBERO-Spatial~\citep{libero} has \NTasks{} tasks. Every task asks the robot to pick up a black bowl and place it on a plate, and the tasks differ only in where the target bowl is (between the plate and the ramekin, on the cookie box, in the top drawer of the wooden cabinet, on the stove, and so on), so the policy has to ground a spatial relation in the image while a second black bowl stands elsewhere on the table. On this suite the instruction is largely redundant with the scene: OpenVLA-OFT still succeeds in \ExtXieInstrLo\% to \ExtXieInstrHi\% of episodes when the instruction is made unreadable~\citep{xie2026language}, so the image carries most of the task information, which is one reason the image steps may matter here more than on suites where the instruction disambiguates the goal. Each task has a fixed file of \NStates{} initial states for evaluation, which gives \NPairs{} episodes per cell. An episode starts with \SettleSteps{} settle steps in which the robot holds still and runs for at most \PolicySteps{} policy steps; it succeeds when the goal predicate of the task holds.

\paragraph{Checkpoints.} We evaluate the released checkpoints unchanged; no model is trained or fine-tuned in this study. OpenVLA-7B was fine-tuned on LIBERO-Spatial with low-rank adaptation (rank \ExtOpenvlaLoraRank), on demonstrations re-rendered at $\RenderRes \times \RenderRes$, rotated by \RotationDeg{} degrees, with no-op actions and the failed demonstrations (\ExtOpenvlaDemosDropped{} of \ExtOpenvlaDemos) removed, using only the third-person camera, and with random-crop augmentation~\citep{openvla,openvlacode}; its authors report \ExtOpenvlaSpatial\% $\pm$ \ExtOpenvlaSpatialSe\% (mean and standard error over \ExtOpenvlaSeeds{} seeds of \NPairs{} episodes). OpenVLA-OFT was fine-tuned from OpenVLA per suite with low-rank adaptation for \ExtOftStepsLo{} to \ExtOftStepsHi{} gradient steps, with parallel decoding, chunks of \Chunk{} actions and an L1 regression loss, and receives the third-person image, a wrist-camera image and the robot state; its authors report \ExtOftSpatial\% on LIBERO-Spatial~\citep{oft}.

\paragraph{Cells and stack.} The study has \NCells{} cells: for each model, the reference, the rerun and \NVariants{} variants, all run between \FirstRunDate{} and \LastRunDate{} \Year{} on one \GpuName{} GPU per cell, with MuJoCo \MujocoVersion{} and greedy decoding for OpenVLA-7B. Appendix~\ref{app:runlog} lists the cells with their dates and host CPUs and the full software stack. The reference cells ran with the original evaluation scripts; the other cells ran with a thin harness that calls the original scripts and replaces only the image function of the step being changed. The host CPU, which can change LIBERO trajectories on its own~\citep{jiang2026bench}, was \CpuMain{} for every harness cell except two: OpenVLA-OFT with the crop step off ran on an \CpuOftCrop{} and OpenVLA-7B rendered at \InputRes{} on an \CpuSevRend; the original scripts do not record the CPU of the reference cells. The OpenVLA-OFT crop-off cell is thus the only OpenVLA-OFT cell whose host CPU is known to differ from the others; a rerun of it on an \CpuMain{} host repeats its outcomes (Appendix~\ref{app:cpucontrol}).
The harness logs one record per episode: task, initial state, outcome, number of steps, wall time, the placement of the cabinet and stove fixtures, hashes of the first frames and the size of the JPEG change on the first input. These records, with the log of the original script for the OpenVLA-OFT reference cell, are the data of every closed-loop table and figure. A pilot on initial states \SplitA, with \PilotCells{} cells of \PilotPairs{} episodes, ran on \PilotDate{} and served only the decision rule of Section~\ref{sec:stats} and the reproducibility check below; it is not pooled with the full cells. The open-loop check ran on \OlDates{} \Year{} on one GPU of the same type.

\paragraph{Same scene in every pair.} Pairing by initial state requires that the paired episodes start from the same scene, and a matching hash of the stored initial state is not enough: LIBERO restores the objects from the stored state but places some fixtures from the environment's own random draw, which depends on the environment's history, so the same stored state can yield different cabinet and stove positions~\citep{issue342}. Before each episode the harness checks the hash of the initial state against the benchmark file, and the hashes match in all pairs of every cell. Beyond that, in all \SceneCells{} harness cells the fixture placement agrees with the other cells of the same model in every pair, and so do the hashes of the first frames, except in the cells rendered at \InputRes, whose frames differ by design while their fixtures agree. All OpenVLA-7B harness cells use environment seed \EnvSeed, and all OpenVLA-OFT cells split the initial states into the same two runs (\SplitA{} and \SplitB) as its reference cell, so every environment goes through the same sequence of resets. The two reference cells ran with the original scripts, which log no fixtures, so for them the match of scenes rests on the reruns, which repeat \SevRerunSame{} and \OftRerunSame{} of \NPairs{} outcomes (Section~\ref{sec:repro}).

\paragraph{Shared runs.} Four cells of this study, the crop-step-off cells and the reference cells of OpenVLA-7B and OpenVLA-OFT on LIBERO-Spatial (\NPairs{} episodes each), are also used in a separate study by the same author, which addresses a different question. That study does not report the results of these cells and refers to this paper for them. No other run, analysis, figure, table or text of this paper appears in that study, and the two studies share only the evaluation harness.
